%!TEX root=../main.tex
\begin{abstract}
Agentic Reinforcement Learning (RL) enables Large Language Models (LLMs) to
perform autonomous decision-making and long-term planning. Unlike standard LLM
post-training, agentic RL workloads are highly \emph{heterogeneous}, combining
compute-intensive prefill phases, bandwidth-bound decoding, and stateful,
CPU-heavy environment simulations. We argue that efficient agentic RL training
requires \emph{disaggregated infrastructure} to leverage specialized, best-fit hardware.
However, naive disaggregation introduces substantial synchronization overhead
and resource underutilization due to the complex dependencies between stages.

We present \SystemName{}, a distributed system designed to maximize throughput for multi-task agentic RL on disaggregated infrastructure. \SystemName{} is built on three core principles: (1) \emph{hardware-affinity workload mapping}, which routes compute-bound and bandwidth-bound tasks to best-fit GPU devices, (2) \emph{fine-grained asynchrony}, which manages execution at the trajectory level to mitigate resource bubbles, and (3) \emph{statefulness-aware computation}, which offloads stateless components (e.g., reward models) to serverless infrastructure for elastic scaling. Our results demonstrate that \SystemName{} effectively improves training throughput and achieves 1.35-2.05\(\times\) end-to-end training time reduction compared to monolithic and synchronous baselines. We also evaluate \SystemName{} by training a hundreds‑of‑billions‑parameter MoE model for Qoder product on an Alibaba cluster with more than 3,000 GPUs, further demonstrating \SystemName{}’s scalability and robustness. The code is available at \url{https://github.com/alibaba/ROLL}.



% Our results demonstrate that \SystemName{} effectively mitigates 
% environment stragglers and achieves significant improvements in training throughput and resource utilization compared to monolithic and synchronous baselines.

% The agentic Reinforcement Learning (RL) emerges as a pivotal technique to enable Large Language Models (LLMs) from passive response to autonomous decision-making and adaptive long-term planning, empowering AI to handle complex and dynamic tasks across domains. RL researchers explore scaling laws in agentic environments and endow LLMs with agentic capability, urging an efficient multi-task agentic RL training system. We assert that efficient multi-task agentic RL necessitates running on disaggregated clusters to leverage specialized hardware. This approach, however, introduces substantial synchronization overhead between stages and can lead to resource underutilization. To address these challenges and fully harness the benefits of disaggregation, we propose \SystemName{}, a system built on three core design principles: (1) map workloads based on hardware affinity; (2) manage execution with fine-grained asynchrony; and (3) characterize computation by statefulness. \SystemName{}’s programming and computation model, as well as its architecture, are designed around these principles. We share using \SystemName{} to train a hundreds‑of‑billions‑parameter MoE model for Qoder product on an Alibaba cluster with more than 3,000 GPUs, and extensive experiments further demonstrate the effectiveness of \SystemName{}.
\end{abstract}


% We also report our experience using \SystemName{} to train a 480B‑A35 Qwen-Coder on an Alibaba cluster with 3{,}072 GPUs, and our extensive empirical analysis demonstrates the effectiveness of \SystemName{}.